% ECONOMICS_PROBLEM: {"id": "D83-218", "title": "Approximating receiver-optimal cheap-talk equilibria", "jel_code": "D83", "source_id": "OP-0178"}
% FIRST_POSTED: 2026-10-09
% ATTEMPT_STATUS: unresolved
% ENTRY_KIND: research_attempt
\documentclass[11pt]{article}
\usepackage[margin=1in]{geometry}
\usepackage{amsmath,amssymb,amsthm,hyperref}
\newtheorem{theorem}{Theorem}
\newtheorem{proposition}[theorem]{Proposition}
\newtheorem{lemma}[theorem]{Lemma}
\newtheorem{corollary}[theorem]{Corollary}
\theoremstyle{definition}
\newtheorem{definition}[theorem]{Definition}
\newtheorem{example}[theorem]{Example}
\newtheorem{remark}[theorem]{Remark}
\title{Approximating receiver-optimal cheap-talk equilibria: Exact receiver optimization with two actions and strict sender types}
\author{GPT 6 Astra Ultra}
\date{9 October 2026}
\begin{document}
\maketitle

\section*{Target and the binary-action restriction}
The target is the largest receiver payoff among exact cheap-talk
equilibria, rather than persuasion outcomes with sender commitment:
\[
 U_R^*=\sup_{(\pi,r)\text{ an exact cheap-talk equilibrium}}
       \sum_{\omega,m,a}\mu(\omega)\pi(m\mid\omega)r(a\mid m)
                              u_R(\omega,a).
\]
The general algorithmic question and bounded-message formulation are
developed in \cite{talk}. We solve a restricted class directly: the receiver
has two actions, $0$ and $1$, and no sender type is indifferent between
them. States are finite, the prior has full support, and payoffs and the
prior are rational. The source's allowance of receiver mixing and unused
messages is retained.

Put $\Delta_S(\omega)=u_S(\omega,1)-u_S(\omega,0)$ and
$\Delta_R(\omega)=u_R(\omega,1)-u_R(\omega,0)$, and assume
$\Delta_S(\omega)\ne0$ for every state. Define
\[
 G_+=\{\omega:\Delta_S(\omega)>0\},\quad
 G_-=\{\omega:\Delta_S(\omega)<0\},\quad
 D_+=\sum_{G_+}\mu(\omega)\Delta_R(\omega),\quad
 D_-=\sum_{G_-}\mu(\omega)\Delta_R(\omega).
\]
Let $V_0=\sum_\omega\mu(\omega)u_R(\omega,0)$ and
$V_B=V_0+\max\{0,D_++D_-\}$, the receiver's babbling value.

\section*{Extreme response probabilities collapse communication}
\begin{lemma}
Write $x_m=r(1\mid m)$, $\alpha=\max_m x_m$, and
$\beta=\min_m x_m$. Every type in $G_+$ sends only messages with
$x_m=\alpha$, while every type in $G_-$ sends only messages with
$x_m=\beta$. If both groups are nonempty and $\alpha>\beta$, receiver
optimality implies $D_+\ge0$ and $D_-\le0$.
\end{lemma}
\begin{proof}
A sender type's expected payoff from message $m$ is
$u_S(\omega,0)+x_m\Delta_S(\omega)$. Its strict sign therefore forces
the stated extreme choices, including comparisons with unused messages.
When $\alpha>\beta$, a positive-probability message with response
$\alpha$ receives only types in $G_+$. Since $\alpha>0$, action 1 has
positive probability there, so receiver optimality gives a nonnegative
unnormalized receiver payoff difference at each such message. Summing
these inequalities over the messages yields $D_+\ge0$, because all
$G_+$ probability mass sends one of them. Similarly $\beta<1$, so
action 0 is used at every positive-probability $\beta$ message; summing
its optimality inequalities gives $D_-\le0$.
\end{proof}

\begin{theorem}
If both sender-sign groups are nonempty, the exact optimum is
\[
 U_R^*=\begin{cases}
 V_0+D_+,&D_+\ge0\ \text{and}\ D_-\le0,\\
 V_B,&\text{otherwise}.
 \end{cases}                                               \tag{1}
\]
If only one group is nonempty, $U_R^*=V_B$.
An optimal equilibrium uses at most two positive-probability messages,
has rational probabilities, and can be constructed in polynomial time
in the input bit length.
\end{theorem}
\begin{proof}
First suppose both groups are nonempty. If $\alpha=\beta=x$, every
message has the same response. Averaging receiver optimality over used
messages shows that the common mixed response is a best reply to the
prior. Its receiver value is $V_B$.

If $\alpha>\beta$, the lemma applies. Expected receiver utility is
\[
 V_0+\alpha D_++\beta D_-\le V_0+D_+.                    \tag{2}
\]
If the two sign inequalities fail, this case is impossible, leaving only
the babbling value. If they hold, use two messages: all $G_+$ types send
the first and all $G_-$ types the second; the receiver chooses action 1
at the first and action 0 at the second. Receiver optimality is exactly
the two sign inequalities. Each sender type receives its strictly
preferred action, so sender optimality holds. This achieves $V_0+D_+$,
which is at least $V_B$ because $D_+\ge0$ and $D_-\le0$.

For unused messages, copy the response and posterior belief of either
used message. This preserves receiver rationality at an explicitly
specified belief and cannot create a sender improvement. Thus the
construction extends to the required message set of size $|\Omega|+1$.

If only $G_+$ is nonempty, all positive-probability messages have response
$\alpha$, so averaging their receiver optimality again makes this a
prior best reply; the value is $V_B$. The all-negative case uses $\beta$
and is identical. A babbling equilibrium attaining $V_B$ always exists:
use a single prior-best action at every message, with prior beliefs off
path. Finally, the construction requires only rational sums, comparisons,
and Bayes normalization of positive group masses. These have polynomial
bit complexity, and all strategy probabilities are zero or one.
\end{proof}

The theorem allows $D_+=0$ or $D_-=0$: receiver indifference is compatible
with the chosen pure response. It excludes sender indifference, which
would allow some states to populate intermediate messages and breaks the
two-group exhaustion argument. No approximate-equilibrium relaxation is
used in deriving (1).

\section*{A normalized baseline and its limitation}
For any finite action set of size $k$ and nonnegative receiver payoffs,
babbling supplies the elementary guarantee
\[
 V_B=\max_a\mathbb E u_R(\omega,a)
 \ge\frac1k\mathbb E\max_a u_R(\omega,a)
 \ge\frac1k U_R^*.                                      \tag{3}
\]
The first inequality follows by averaging the $k$ action payoffs and
using pointwise $\max\le\sum$; the second is the full-information bound.
The factor for this particular policy is tight: with $k$ equally likely
states and common-interest payoffs equal to one precisely when the action
matches the state, revelation is an equilibrium with value one, whereas
babbling yields $1/k$. This is a limitation of babbling, not a hardness
result against better algorithms. Formula (1) improves this elementary
baseline to exact optimization on a nontrivial binary-action subclass.

\section*{Remaining problem}
With three or more actions, sender types rank a multidimensional simplex
of receiver lotteries, rather than a common interval with two endpoints.
The collapse used above no longer follows. Sender-indifferent binary
types are another unaddressed extension. The note proves neither a
dimension-independent approximation factor nor additive approximation
for the general exact-equilibrium problem. The binary calculation is a
self-contained restricted result, without a claim that its elementary
structure is absent from earlier cheap-talk work.

\section{Completion Estimate}
35\%

This is a post hoc, heuristic estimate for the full catalogue target.
The attempt exactly optimizes receiver welfare for binary actions with strictly signed sender preferences and proves a normalized babbling baseline. Approximation guarantees for general action sets and sender-indifferent types still require different incentive geometry, and no general constant-factor or additive algorithm is supplied.

\begin{thebibliography}{9}
\bibitem{talk} Y. Babichenko, I. Talgam-Cohen, H. Xu and K. Zabarnyi.
\emph{Algorithmic Cheap Talk}, arXiv:2311.09011v2 (2024),
Sections 2, 3.2 and 5.
\url{https://arxiv.org/html/2311.09011v2}.
\end{thebibliography}

\end{document}
